\section{Introduction}

Recommender systems have evolved from early accuracy-focused algorithms 
to complex industrial systems across e-commerce, media, and social platforms.
Despite these advances, a tension exists between short-term business-driven goals and long-term user wellbeing.
Most current systems rely on easily measurable signals like clicks and watch time, 
which maximize immediate engagement but often do not explicitly model the user's cognitive state.

\subsection*{The Challenge of Engagement Bias}

Standard sequential and deep learning models treat user histories as ordered sequences to predict the next likely interaction.
While effective at capturing recency effects, these models implicitly assume that maximizing immediate clicks corresponds to maximizing satisfaction.
This approach neglects psychological factors such as \textbf{fatigue} and \textbf{diminishing utility}, 
where repeated exposure to similar content reduces its value and this leads to platform disengagement.

\subsection*{Proposed Solution: DS-HCRS}

To address these limitations, this work proposes a framework called \textbf{DS-HCRS} 
which goes beyond a single static strategy by incorporating:

\begin{itemize}
    \item \textbf{Dual-System Architecture:} Balancing a fast "impulsive" model for relevance with a reflective "long-term" optimizer.
    \item \textbf{Dynamics of Preference:} Estimating changes in user preference using recurrent models.
    \item \textbf{Cognitive State Modeling:} Estimating user fatigue and discontentment using Reinforcement Learning objectives.
\end{itemize}

These three systems alone allows the algorithm to implement more complex systems like 
updating user preference (embeddings) on the go, and cold start handling for new users and items.
By framing recommendation as a cognitive process rather than just a prediction problem,
we aim to develop systems that are sustainable, adaptive, and better aligned with long-term user experience.
